\begingroup
\renewcommand{\clearpageforsection}{}
\exercisesheader{}
\endgroup


% 35

\eoce{\qt{Berhenti merokok\label{quitters_chisq_independence}} Apakah keikutsertaan dalam 
kelompok dukungan memengaruhi kemampuan seseorang untuk berhenti merokok? Dinas 
kesehatan sebuah kabupaten mengikutsertakan 300 perokok dalam eksperimen acak. Sebanyak 150 
peserta ditempatkan dalam kelompok yang menggunakan koyo nikotin dan 
bertemu setiap minggu dengan kelompok dukungan; 150 peserta lainnya menerima koyo dan 
tidak bertemu dengan kelompok dukungan. Pada akhir penelitian, 40 
peserta dalam kelompok koyo plus dukungan telah berhenti merokok, sedangkan 
hanya 30 perokok yang berhenti dalam kelompok lainnya.
\begin{parts}
\item Buatlah tabel dua arah yang menyajikan hasil penelitian ini.
\item Jawablah setiap pertanyaan berikut di bawah hipotesis nol 
bahwa keikutsertaan dalam kelompok dukungan tidak memengaruhi kemampuan 
seseorang untuk berhenti merokok, dan nyatakan apakah nilai harapannya 
lebih tinggi atau lebih rendah daripada nilai teramati.
\begin{subparts}
\item Berapa banyak subjek dalam kelompok ``koyo + dukungan" yang Anda 
harapkan berhenti merokok?
\item Berapa banyak subjek dalam kelompok ``hanya koyo" yang Anda harapkan 
tidak berhenti merokok?
\end{subparts}
\end{parts}
}{}

% 36

\eoce{\qt{Pemindaian seluruh tubuh, Bagian II\label{full_body_scan_chisq_indep}} Tabel 
di bawah merangkum gugus data yang pertama kali kita jumpai dalam 
Latihan~\ref{full_body_scan_HT_Error} mengenai pandangan terhadap pemindaian seluruh 
tubuh dan afiliasi politik. Perbedaan pada setiap kelompok politik 
mungkin terjadi karena kebetulan. Selesaikan perhitungan berikut di bawah 
hipotesis nol bahwa afiliasi partai seseorang saling bebas dengan 
dukungannya terhadap pemindaian seluruh tubuh. Sebelum melanjutkan, mungkin berguna untuk 
terlebih dahulu menambahkan satu kolom tambahan bagi total baris sebelum melakukan 
perhitungan.
\begin{center}
\begin{tabular}{ll  cc c} 
            &   & \multicolumn{3}{c}{\textit{Afiliasi Partai}} \\
\cline{3-5}
                                &           & Republik & Demokrat & Independen   \\
\cline{2-5}
\multirow{3}{*}{\textit{Jawaban}}& Sebaiknya    & 264        & 299      & 351 \\
                                & Sebaiknya tidak& 38         & 55       & 77 \\
                                & Tidak tahu/Tidak menjawab & 16 & 15    & 22 \\
\cline{2-5}
                                & Total      & 318       & 369      & 450
\end{tabular}
\end{center}
\begin{parts}
\item Berapa banyak pemilih Republik yang Anda harapkan tidak mendukung penggunaan 
pemindaian seluruh tubuh?
\item Berapa banyak pemilih Demokrat yang Anda harapkan mendukung penggunaan pemindaian seluruh 
tubuh?
\item Berapa banyak pemilih Independen yang Anda harapkan tidak tahu atau tidak menjawab?
\end{parts}
}{}

% 37

\eoce{\qt{Pengeboran lepas pantai, Bagian III\label{offshore_drilling_chisq_indep}} 
Tabel di bawah merangkum gugus data yang pertama kali kita jumpai dalam 
Latihan~\ref{offshore_drill_edu_dontknow_HT}, yang menelaah 
tanggapan sampel acak lulusan perguruan tinggi dan bukan lulusan mengenai 
pengeboran minyak. Lakukan uji khi-kuadrat pada data ini untuk 
memeriksa apakah terdapat perbedaan yang signifikan secara statistik antara 
tanggapan lulusan perguruan tinggi dan bukan lulusan.
\begin{center}
\begin{tabular}{l c c}
			& \multicolumn{2}{c}{\textit{Lulusan Perguruan Tinggi}} \\
\cline{2-3}
			& Ya		& Tidak				\\
\cline{1-3}
Mendukung		& 154		& 132			\\
Menentang		& 180		& 126			\\
Tidak tahu	& 104		& 131			\\
\cline{1-3}
 Total		& 438		& 389		
\end{tabular}
\end{center}
}{}

% 38

\eoce{\qt{Cacing parasit\label{parasitic_worm_chisq}}
Filariasis limfatik adalah penyakit yang disebabkan oleh cacing parasit.
Komplikasi penyakit ini dapat menyebabkan pembengkakan ekstrem
dan komplikasi lainnya.
Di sini kita mempertimbangkan hasil eksperimen acak
yang membandingkan tiga
pilihan pengobatan berbeda untuk membersihkan tubuh seseorang dari
parasit ini, yang sedang diupayakan agar sepenuhnya dimusnahkan.
Hasil penelitian untuk tahun kedua
diberikan di bawah ini:\footfullcite{King_Suamani_2018}
\begin{center}
\begin{tabular}{l cc}
  \hline
  & Bersih pada Tahun ke-2 & Belum Bersih pada Tahun ke-2 \\ 
  \hline
  Tiga obat & 52 & 2 \\ 
  Dua obat & 31 & 24 \\ 
  Dua obat setiap tahun & 42 & 14 \\ 
  \hline
\end{tabular}
\end{center}
\begin{parts}
\item\label{parasitic_worm_chisq_hyp}
    Susunlah hipotesis untuk mengevaluasi
    apakah terdapat perbedaan dalam
    kinerja berbagai pengobatan,
    dan periksa pula syarat-syaratnya.
\item
    Perangkat lunak statistik digunakan untuk menjalankan
    uji khi-kuadrat, yang menghasilkan:
    \begin{align*}
    &X^2 = 23.7
    &&df = 2
    &&\text{nilai-p} = \text{7.2e-6}
    \end{align*}
    Gunakan hasil ini untuk mengevaluasi hipotesis
    dari bagian~(\ref{parasitic_worm_chisq_hyp}),
    dan berikan kesimpulan
    dalam konteks masalah ini.
\end{parts}
}{}
